\documentclass[../../main.tex]{subfiles}

\begin{document}

\chapter{Analisis Algoritma Block Cipher: DES}
\label{chap:des}

\begin{learningoutcome}
Setelah mempelajari bab ini, mahasiswa diharapkan mampu:
\begin{itemize}
    \item Menjelaskan sejarah, struktur, dan mekanisme kerja algoritma Data Encryption Standard (DES) (Sub-CPMK 1.2).
    \item Menganalisis proses pembangkitan kunci internal (\textit{key scheduling}) pada DES (Sub-CPMK 1.2).
    \item Mengevaluasi keamanan DES dan memahami risiko serangan \textit{brute force} serta \textit{meet-in-the-middle} pada Double DES (Sub-CPMK 1.2).
    \item Menjelaskan mekanisme kerja Triple DES (3DES) sebagai solusi peningkatan keamanan DES (Sub-CPMK 1.2).
\end{itemize}
\end{learningoutcome}

\subfile{section-01}
\subfile{section-02}
\subfile{section-03}
\section{Contoh Terhitung}

\begin{example}
\textbf{Satu putaran jaringan Feistel.} Blok mainan 8 bit dibagi dua bagian 4 bit; fungsi putaran $f(R,K)=(R\oplus K)$ digeser siklik ke kiri 1 bit. Mula-mula $L_0=1010$, $R_0=0011$, dengan aturan
\[
L_i=R_{i-1},\qquad R_i=L_{i-1}\oplus f(R_{i-1},K_i).
\]
Putaran 1, $K_1=1100$: $f(0011,1100)=\text{geser-kiri}(0011\oplus1100)=\text{geser-kiri}(1111)=1111$, sehingga $L_1=0011$ dan $R_1=1010\oplus1111=0101$.\\
Putaran 2, $K_2=1010$: $f(0101,1010)=\text{geser-kiri}(1111)=1111$, sehingga $L_2=0101$ dan $R_2=0011\oplus1111=1100$.\\
Melanjutkan dengan $K_3=0110$ dan $K_4=1001$:
\[
\begin{array}{c|cc}
i & L_i & R_i\\
\hline
0 & 1010 & 0011\\
1 & 0011 & 0101\\
2 & 0101 & 1100\\
3 & 1100 & 0000\\
4 & 0000 & 1111
\end{array}
\]
Dekripsi memakai aturan yang sama dengan urutan kunci terbalik $(K_4,K_3,K_2,K_1)$ dan menukar peran $L,R$, tanpa perlu fungsi $f^{-1}$ yang terpisah.
\end{example}

\begin{example}
\textbf{Longsoran dan pembangkitan subkunci DES.}
(i) \textit{Efek longsoran.} Ubah satu bit plainteks di atas, $L_0:1010\to0010$. Setelah putaran 1, 2, 3, dan 4, jumlah bit keluaran $(L_i,R_i)$ yang berbeda dari semula berturut-turut adalah $1,2,3,3$ bit (dari 8 bit). Satu bit masukan makin menyebar seiring bertambahnya putaran; pada DES penuh (16 putaran, blok 64 bit) perubahan 1 bit menyebar ke sekitar setengah bit cipherteks.
(ii) \textit{Pembangkitan subkunci.} Kunci efektif 56 bit dipecah menjadi $C_0,D_0$ (masing-masing 28 bit). Untuk $C_0=1111000011110000111100001111$:
\[
\text{putaran 1, geser-kiri 1}\to C_1=1110000111100001111000011111,
\]
\[
\text{putaran 2, geser-kiri 1 lagi}\to C_2=1100001111000011110000111111,
\]
sesuai jadwal pergeseran putaran 1 dan 2 (masing-masing 1 bit). Hasil $C_i,D_i$ lalu dipilih 48 bit oleh PC-2 menjadi $K_i$.
(iii) \textit{Substitusi S-box.} S-box $S_1$ memetakan 6 bit ke 4 bit: bit terluar menentukan baris, empat bit tengah menentukan kolom. Untuk masukan $011011$: baris $01_2=1$, kolom $1101_2=13$, dan $S_1(1,13)=5$, jadi keluarannya $0101$.
\end{example}


\begin{obeactivity}
\textbf{Aktivitas 7.1: Analisis Kompleksitas Kunci}
1. Hitung berapa lama waktu yang dibutuhkan untuk melakukan brute force terhadap DES (56 bit) jika sebuah komputer dapat mencoba $10^9$ kunci per detik.
2. Bandingkan hasil tersebut dengan waktu yang dibutuhkan untuk menyerang AES-128 ($2^{128}$ kemungkinan).
3. Diskusikan mengapa peningkatan panjang kunci secara eksponensial jauh lebih efektif daripada sekadar mengulang algoritma yang sama (seperti pada Double DES).
\end{obeactivity}


\begin{obereflection}
\textbf{Latihan 7.1: Konsep}
1. Mengapa Triple DES menggunakan urutan Enkripsi-Dekripsi-Enkripsi (EDE) alih-alih Enkripsi-Enkripsi-Enkripsi (EEE)?
2. Jelaskan mekanisme kerja serangan \textit{Meet-in-the-Middle} pada Double DES secara sederhana.
3. Apa dampak dari penggunaan kunci lemah (\textit{weak keys}) pada algoritma DES?
4. Mengapa bit paritas pada kunci DES tidak menambah kekuatan kunci? Berapa panjang kunci efektifnya?
5. Jelaskan mengapa permutasi awal (IP) dan inversinya tidak menambah keamanan DES sama sekali.
6. Apa fungsi tumpang tindih bit pada tabel ekspansi $E$ terhadap penyebaran pengaruh satu bit?
7. Mengapa kriptanalisis linier Matsui, meskipun secara teori memerlukan lebih sedikit operasi daripada pencarian menyeluruh, tidak pernah dipakai untuk memecahkan DES dalam praktik?
8. Mengapa batas $2^{32}$ blok lebih berbahaya bagi 3DES daripada bagi AES-128?
\end{obereflection}

\begin{obereflection}
\textbf{Latihan 7.2: Hitungan dengan Tabel Bit}
Kerjakan dengan tabel pada Tabel~\ref{tab:ip-des}, Tabel~\ref{tab:ep-des},
Tabel~\ref{tab:sbox-des}, dan Tabel~\ref{tab:pc-des}.

1. Tentukan $C_0$ dan $D_0$ untuk kunci $\code{0000000000000000}$ (semua bit nol).
   Apakah kunci ini termasuk kunci lemah? Jelaskan.
2. Tentukan keluaran $S_1$ untuk masukan $\code{001011}$ dan $\code{111010}$.
3. Pada tabel ekspansi $E$, sebutkan posisi bit masukan 32 bit yang muncul
   \emph{dua kali} pada keluaran 48 bit.
4. Pada Kunci DES, posisi berapa saja yang merupakan bit paritas? Berapa banyak?
5. Berapa banyak subkunci yang dibangkitkan DES dan berapa panjang masing-masing?
6. Tentukan $\text{IP}(\code{8000000000000000})$, yaitu permutasi awal atas blok
   yang hanya bit pertamanya bernilai 1. Nyatakan jawabannya dalam heksadesimal.
7. Hitung waktu rata-rata pemecahan DES pada laju $10^{10}$ kunci per detik.
   Nyatakan jawabannya dalam hari.
8. Berapa entri yang harus disimpan penyerang pada serangan
   \textit{meet-in-the-middle} terhadap Double DES, dan berapa operasi yang
   dibutuhkan seluruhnya?
\end{obereflection}

\section{Rangkuman}
\begin{itemize}
    \item DES adalah \textit{block cipher} blok 64 bit dengan kunci eksternal 64 bit yang hanya 56 bit dipakai untuk enkripsi (8 bit sisanya paritas), memakai jaringan Feistel 16 putaran.
    \item Alur enkripsi: permutasi awal (IP), 16 putaran \textit{enciphering}, lalu invers permutasi awal ($\text{IP}^{-1}$).
    \item Fungsi putaran $f$ memperluas blok kanan 32 bit menjadi 48 bit, di-XOR dengan subkunci, disubstitusi melalui delapan S-box (6 bit $\to$ 4 bit), lalu dipermutasi oleh $P$.
    \item Pembangkitan kunci: PC-1 membuang bit paritas dan membagi kunci menjadi $C_0,D_0$ (28 bit), pergeseran siklik per putaran, dan PC-2 memilih 48 bit menjadi $K_i$.
    \item Sifat Feistel yang reversibel membuat dekripsi memakai algoritma yang sama dengan urutan kunci terbalik $K_{16},\dots,K_1$, tanpa fungsi dekripsi terpisah.
    \item Keamanan DES terbatas pada ruang kunci $2^{56}$ yang kini terlalu kecil untuk pencarian menyeluruh dan adanya \textit{weak keys}; efek longsorannya kuat sehingga perubahan satu bit plainteks atau kunci menyebar ke sekitar setengah bit cipherteks.
    \item Double DES ($2^{112}$) tumbang oleh \textit{meet-in-the-middle} sekitar $2^{57}$; Triple DES EDE (2 kunci 112 bit, 3 kunci 168 bit) memperpanjang umur DES.
\end{itemize}


\begin{obeassessment}
\textbf{Kuis Bab 7}
1. Jelaskan perbedaan antara kunci eksternal dan kunci putaran (\textit{round key}) pada DES, serta sebutkan panjang masing-masing.
2. Mengapa DES dianggap tidak aman untuk standar keamanan modern, meskipun memiliki struktur Jaringan Feistel yang kompleks?
3. Bagaimana Triple DES meningkatkan keamanan dibandingkan dengan Single DES?
4. Pada pembangkitan subkunci DES, mengapa PC-2 hanya memilih 48 dari 56 bit yang tersedia, dan apa akibatnya bagi pemakaian bit kunci?
5. Diberikan masukan $S_3$ sebesar $\code{011110}$. Tentukan nomor baris, nomor kolom, dan keluaran 4 bitnya dengan membaca Tabel~\ref{tab:sbox-des}.
6. Sebuah sistem memakai DES-CBC. Jelaskan mengapa memakai ulang IV yang sama untuk dua pesan berbeda melemahkan keamanan, dan bagaimana hubungannya dengan mode ECB.
7. Buktikan bahwa pada kunci $\code{FEFEFEFEFEFEFEFE}$, seluruh 16 subkunci DES bernilai sama.
8. Seorang insinyur mengusulkan mengganti 3DES tiga kunci dengan Double DES karena keduanya sama-sama "dua kali DES". Bantah usulan itu dengan membandingkan kompleksitas serangan nyata pada keduanya.
\end{obeassessment}
\begin{competencychecklist}
\begin{tabularx}{\linewidth}{|X|c|}
\hline
Indikator Kompetensi & Tercapai \\
\hline
Saya dapat menjelaskan struktur global algoritma DES & [ ] \\
\hline
Saya dapat menjelaskan sejarah DES dari Lucifer sampai FIPS PUB 46 & [ ] \\
\hline
Saya dapat menguraikan proses pembangkitan kunci internal DES & [ ] \\
\hline
Saya dapat membaca tabel IP, $\text{IP}^{-1}$, $E$, dan $P$ pada blok 64 bit & [ ] \\
\hline
Saya dapat mensubstitusi satu kelompok 6 bit melalui S-box DES & [ ] \\
\hline
Saya dapat mengerjakan satu putaran DES lengkap di atas kertas & [ ] \\
\hline
Saya dapat menganalisis kelemahan Double DES terhadap Meet-in-the-Middle attack & [ ] \\
\hline
Saya dapat membedakan varian Triple DES (2 kunci vs 3 kunci) & [ ] \\
\hline
Saya dapat menyebutkan nilai kunci lemah dan menjelaskan mengapa ia lemah & [ ] \\
\hline
Saya dapat menjelaskan konsep efek longsoran (\textit{avalanche effect}) & [ ] \\
\hline
Saya dapat menjelaskan mengapa kriptanalisis linier tidak dipakai dalam praktik & [ ] \\
\hline
Saya dapat menghitung waktu pemecahan DES pada laju pencarian tertentu & [ ] \\
\hline
Saya dapat menjelaskan alasan 3DES ditinggalkan setelah AES dibakukan & [ ] \\
\hline
\end{tabularx}
\end{competencychecklist}


\end{document}
